\subsection{Filtered ring associated with a graded ring}

For every \(n \in \mathbf{Z}\) , we set \(\mathrm{H}_{\geqslant n} = \sum_{i \geqslant n} \mathrm{H}_i\) . We have \(\mathrm{H} = \mathrm{H}_{\geqslant 0}\) ; the \(\mathrm{H}_{\geqslant n}\) are graded ideals of \(\mathrm{H}\) . Let S be the multiplicative subset \(1 + \mathrm{H}_{\geqslant 1}\) formed by the elements of H whose component of degree 0 is equal to 1, and consider the ring of fractions \(\mathrm{S}^{-1}\mathrm{H}\) . Identify H with a subring of its completion \(\hat{\mathrm{H}} = \prod_{n} \mathrm{H}_{n}\) (III, § 2, no. 12, example 1); since the elements of S are invertible in \(\hat{\mathrm{H}}\) (III, § 2, no. 13, lemma 3), \(\mathrm{S}^{-1}\mathrm{H}\) is identified with a subring of \(\hat{\mathrm{H}}\) containing H. For \(s \in \mathrm{S}\) and \(h \in \mathrm{H}_{\geqslant n}\) , the element \(s^{-1}h - h\) of \(\hat{\mathrm{H}}\) belongs to \(\prod_{i \geqslant n} \mathrm{H}_i\) ; consequently we have \(\mathrm{S}^{-1}\mathrm{H}_{\geqslant n} = (\mathrm{S}^{-1}\mathrm{H}) \cap \prod_{i \geqslant n} \mathrm{H}_i\) .

We deduce from this:

PROPOSITION 1. --- a) The ideals S \(^{-1}\) H \(_{≥n}\) form an exhaustive and separated filtration of the ring S \(^{-1}\) H.

\begin{enumerate}
\def\labelenumi{\alph{enumi})}
\setcounter{enumi}{1}
\tightlist
\item
  The canonical homomorphism from H to \(S^{-1}H\) induces for each n an isomorphism \(u_{n}\) of \(H_{n}\) onto \(S^{-1}H_{\geqslant n}/S^{-1}H_{\geqslant n+1}\) ; the \(u_{n}\) are the homogeneous components of an isomorphism of graded rings from H onto the graded ring associated with \(S^{-1}H\) , filtered by the \(S^{-1}H_{\geqslant n}\) .
\end{enumerate}

Remarks. --- 1) An element \(h/s\) of \(\mathrm{S}^{-1}\mathrm{H}\) with \(h\in\mathrm{H}, s\in\mathrm{S}\) , is invertible if and only if the component of degree 0 of \(h\) is invertible in \(\mathrm{H}_{0}\) . Consequently, if the ring \(\mathrm{H}_{0}\) is local, the ring \(\mathrm{S}^{-1}\mathrm{H}\) is local and the canonical injection \(\mathrm{H}_{0}\to\mathrm{S}^{-1}\mathrm{H}\) induces an isomorphism on the residue fields.

\begin{enumerate}
\def\labelenumi{\arabic{enumi})}
\setcounter{enumi}{1}
\tightlist
\item
  Suppose H is generated by \(H_{0}\) and \(H_{1}\) ; then for every n, we have \(H_{n+1} = H_{1}.H_{n}\) , hence \(H_{\geqslant n+1} = H_{1}.H_{\geqslant n}\) and \(S^{-1}H_{\geqslant n+1} = H_{1}.S^{-1}H_{\geqslant n}\) . Consequently, the filtration \((S^{-1}H_{\geqslant n})\) of \(S^{-1}H\) is the filtration \(S^{-1}H_{\geqslant 1}\) -adic.
\end{enumerate}

Examples. --- 1) * Let p be a graded prime ideal of C{[}X₀, \ldots, Xₙ{]} different from the ideal generated by the Xᵢ; let V be the algebraic subvariety of Pⁿ(C) defined by p and C the algebraic subvariety of Cⁿ⁺¹ defined by p. Then C is the cone over V, H = C{[}X₀, \ldots, Xₙ{]}/p is the affine algebra of C and S⁻¹H is the local ring of the cone C at its vertex. *

\begin{enumerate}
\def\labelenumi{\arabic{enumi})}
\setcounter{enumi}{1}
\tightlist
\item
  Let A be a local ring and a an ideal of A distinct from A. Then H = \(\bigoplus_{n} \mathfrak{a}^{n}/\mathfrak{a}^{n+1}\)
\end{enumerate}

is a graded ring such that \(H_{0} = A/\mathfrak{a}\) is local; it is generated by \(H_{0}\) and \(H_{1}\) . The ring \(S^{-1}H\) is therefore local and the filtration ( \(S^{-1}H_{\geq n}\) ) is the \(S^{-1}H_{\geq 1}\) -adic filtration. Note that in general the rings A and \(S^{-1}H\) are not isomorphic. * In particular, an algebraic variety is not in general locally isomorphic, in a neighborhood of a point, to its tangent cone at that point. *

